\documentclass[11pt,a4paper]{article}

\usepackage{amsmath,amssymb,amsthm}
\usepackage{mathtools}
\usepackage[margin=1in]{geometry}
\usepackage{stmaryrd}

% Custom inference rule macro (replaces proof.sty)
\newcommand{\infer}[3][]{%
  \genfrac{}{}{0.4pt}{0}{#3}{#2}\;\textsc{\footnotesize #1}%
}

\newtheorem{definition}{Definition}[section]
\newtheorem{remark}{Remark}[section]
\newtheorem{theorem}{Theorem}[section]
\newtheorem{lemma}[theorem]{Lemma}
\newtheorem{example}[theorem]{Example}

\newcommand{\W}{\mathsf{W}}
\newcommand{\Type}{\mathsf{Type}}
\newcommand{\ctx}{\;\mathsf{ctx}}
\newcommand{\type}{\;\mathsf{type}}
\newcommand{\at}[1]{\,@\,#1}
\newcommand{\one}{\mathbf{1}}
\newcommand{\zero}{\mathbf{0}}

\title{ProbTT: Type Theory with Primitive Weights\\[0.5em]
\large Minimal Axiomatization}

\author{Probabilistic Foundations Project}
\date{\today}

\begin{document}

\maketitle

\begin{abstract}
We define a type theory where judgments carry \emph{weights} from a primitive weight monoid $\W$. The monoid has multiplication, unit, zero, and order---but NO addition or subtraction. This is the minimal structure needed for combining uncertainties. When $\W = \{\zero, \one\}$, we recover Martin-L\"of Type Theory. When $\W = [0,1]$, we get probabilistic typing. With Markov category structure (copy, delete, disintegration), we get conditioning and Bayes' rule. The weight structure is axiomatized, not constructed from real numbers.
\end{abstract}

\tableofcontents

%=============================================================================
\section{The Weight Monoid}
%=============================================================================

\subsection{Minimal Structure}

We do NOT assume numbers. We axiomatize the minimal structure needed for weights.

\begin{definition}[Weight Monoid]
$\W$ is a type with:
\begin{itemize}
    \item Elements: $\one : \W$ (certainty), $\zero : \W$ (impossibility)
    \item Operation: $(\cdot) : \W \to \W \to \W$ (combination)
    \item Order: $(\leq) : \W \to \W \to \Type$ (sub-weight)
\end{itemize}
\end{definition}

\begin{definition}[Weight Axioms]
The following equations and laws hold:
\begin{align}
w \cdot \one &= w \tag{unit-right} \\
\one \cdot w &= w \tag{unit-left} \\
w \cdot \zero &= \zero \tag{annihilate-right} \\
\zero \cdot w &= \zero \tag{annihilate-left} \\
(u \cdot v) \cdot w &= u \cdot (v \cdot w) \tag{assoc} \\
u \cdot v &= v \cdot u \tag{comm}
\end{align}
\end{definition}

\begin{definition}[Order Axioms]
\begin{align}
w &\leq w \tag{refl} \\
u \leq v \land v \leq w &\Rightarrow u \leq w \tag{trans} \\
u \leq v \land v \leq u &\Rightarrow u = v \tag{antisym} \\
\zero &\leq w \tag{zero-least} \\
w &\leq \one \tag{one-greatest} \\
u \leq v &\Rightarrow u \cdot w \leq v \cdot w \tag{mono}
\end{align}
\end{definition}

\begin{remark}[What We DON'T Have]
\begin{itemize}
    \item NO addition $+$
    \item NO subtraction $-$
    \item NO complement $\one - w$
    \item NO intermediate values (yet)
\end{itemize}
This is a \textbf{bounded ordered commutative monoid}, not a semiring.
\end{remark}

\subsection{Examples}

\begin{example}[Boolean Weights]
$\W = \{\zero, \one\}$ with $\one \cdot \one = \one$, else $\zero$.

This gives classical logic: certain or impossible, nothing between.
\end{example}

\begin{example}[Probability Weights]
$\W = [0,1] \subset \mathbb{R}$ with ordinary multiplication.

This gives probabilistic reasoning.
\end{example}

\begin{example}[Discrete Weights]
$\W = \{0, \frac{1}{n}, \frac{2}{n}, \ldots, 1\}$ for fixed $n$.

Finite, decidable, approximates probability.
\end{example}

\begin{example}[Tropical Weights]
$\W = [0, \infty]$ with $u \cdot v = u + v$ (tropical multiplication).

This gives cost/distance semantics instead of probability.
\end{example}

%=============================================================================
\section{Judgments}
%=============================================================================

\begin{definition}[Judgment Forms]
\begin{enumerate}
    \item $\Gamma \ctx$ \hfill (valid context)
    \item $\Gamma \vdash A \type$ \hfill (valid type)
    \item $\Gamma \vdash a : A \at{w}$ \hfill (weighted typing)
    \item $\Gamma \vdash a \equiv b : A \at{w}$ \hfill (weighted equality)
\end{enumerate}
where $w$ is a closed term of type $\W$.
\end{definition}

\begin{remark}[Interpretation]
$\Gamma \vdash a : A \at{w}$ means: term $a$ inhabits type $A$ with weight $w$.
\begin{itemize}
    \item $w = \one$: $a$ definitely inhabits $A$
    \item $w = \zero$: $a$ does not inhabit $A$ (vacuously true)
    \item $\zero < w < \one$: $a$ partially/probably inhabits $A$
\end{itemize}
\end{remark}

%=============================================================================
\section{Structural Rules}
%=============================================================================

\subsection{Contexts}

\begin{gather}
\infer[\textsc{Ctx-Empty}]
  {\cdot \ctx}
  {}
\\[1em]
\infer[\textsc{Ctx-Ext}]
  {\Gamma, x : A \ctx}
  {\Gamma \ctx \quad \Gamma \vdash A \type}
\end{gather}

Contexts have no weights---they list available assumptions.

\subsection{Variables}

\begin{gather}
\infer[\textsc{Var}]
  {\Gamma, x : A, \Delta \vdash x : A \at{\one}}
  {}
\end{gather}

Variables have weight $\one$: if $x : A$ is assumed, then $x$ definitely has type $A$.

\subsection{Weakening}

\begin{gather}
\infer[\textsc{Weaken}]
  {\Gamma \vdash a : A \at{v}}
  {\Gamma \vdash a : A \at{w} \quad v \leq w}
\end{gather}

If $a : A$ with weight $w$, then $a : A$ with any smaller weight $v$.

\subsection{Substitution}

\begin{gather}
\infer[\textsc{Subst}]
  {\Gamma \vdash b[a/x] : B[a/x] \at{w \cdot v}}
  {\Gamma \vdash a : A \at{w} \quad \Gamma, x : A \vdash b : B \at{v}}
\end{gather}

Substituting a weighted term multiplies weights.

%=============================================================================
\section{Weight Type}
%=============================================================================

\begin{gather}
\infer[\textsc{W-Form}]
  {\Gamma \vdash \W \type}
  {\Gamma \ctx}
\\[1em]
\infer[\textsc{W-Zero}]
  {\Gamma \vdash \zero : \W \at{\one}}
  {\Gamma \ctx}
\qquad
\infer[\textsc{W-One}]
  {\Gamma \vdash \one : \W \at{\one}}
  {\Gamma \ctx}
\\[1em]
\infer[\textsc{W-Mul}]
  {\Gamma \vdash u \cdot v : \W \at{\one}}
  {\Gamma \vdash u : \W \at{\one} \quad \Gamma \vdash v : \W \at{\one}}
\end{gather}

Weight terms have weight $\one$: the weights themselves are certain.

%=============================================================================
\section{Function Types}
%=============================================================================

\subsection{Formation}

\begin{gather}
\infer[\textsc{Fun-Form}]
  {\Gamma \vdash (x : A) \to B \type}
  {\Gamma \vdash A \type \quad \Gamma, x : A \vdash B \type}
\end{gather}

Function types are certain---no weight on type formation.

\subsection{Introduction}

\begin{gather}
\infer[\textsc{Fun-Intro}]
  {\Gamma \vdash \lambda x. b : (x : A) \to B \at{w}}
  {\Gamma, x : A \vdash b : B \at{w}}
\end{gather}

A function has weight $w$ if its body has weight $w$ (uniformly in $x$).

\subsection{Elimination}

\begin{gather}
\infer[\textsc{Fun-Elim}]
  {\Gamma \vdash f\,a : B[a/x] \at{w \cdot v}}
  {\Gamma \vdash f : (x : A) \to B \at{w} \quad \Gamma \vdash a : A \at{v}}
\end{gather}

\textbf{Key rule}: weights multiply. Uncertain function applied to uncertain argument gives product weight.

\subsection{Computation}

\begin{gather}
\infer[\textsc{Fun-$\beta$}]
  {\Gamma \vdash (\lambda x. b)\,a \equiv b[a/x] : B[a/x] \at{w}}
  {\Gamma, x : A \vdash b : B \at{w} \quad \Gamma \vdash a : A \at{\one}}
\end{gather}

$\beta$-reduction preserves weight (when argument is certain).

%=============================================================================
\section{Pair Types}
%=============================================================================

\subsection{Formation}

\begin{gather}
\infer[\textsc{Pair-Form}]
  {\Gamma \vdash \Sigma(x : A). B \type}
  {\Gamma \vdash A \type \quad \Gamma, x : A \vdash B \type}
\end{gather}

\subsection{Introduction}

\begin{gather}
\infer[\textsc{Pair-Intro}]
  {\Gamma \vdash (a, b) : \Sigma(x : A). B \at{w \cdot v}}
  {\Gamma \vdash a : A \at{w} \quad \Gamma \vdash b : B[a/x] \at{v}}
\end{gather}

Pair weight = product of component weights. (Joint probability.)

\subsection{Elimination}

\begin{gather}
\infer[\textsc{Pair-Fst}]
  {\Gamma \vdash \pi_1\,t : A \at{w}}
  {\Gamma \vdash t : \Sigma(x : A). B \at{w}}
\\[1em]
\infer[\textsc{Pair-Snd}]
  {\Gamma \vdash \pi_2\,t : B[\pi_1 t / x] \at{w}}
  {\Gamma \vdash t : \Sigma(x : A). B \at{w}}
\end{gather}

Projections preserve weight.

%=============================================================================
\section{Identity Types}
%=============================================================================

\subsection{Formation}

\begin{gather}
\infer[\textsc{Id-Form}]
  {\Gamma \vdash \mathsf{Id}_A(a, b) \type}
  {\Gamma \vdash a : A \at{\one} \quad \Gamma \vdash b : A \at{\one}}
\end{gather}

Identity types require certain endpoints.

\subsection{Introduction}

\begin{gather}
\infer[\textsc{Id-Intro}]
  {\Gamma \vdash \mathsf{refl} : \mathsf{Id}_A(a, a) \at{w}}
  {\Gamma \vdash a : A \at{w}}
\end{gather}

Reflexivity has the same weight as the term.

\subsection{Elimination}

\begin{gather}
\infer[\textsc{Id-Elim}]
  {\Gamma \vdash J(d, e) : P[b, e] \at{w \cdot v}}
  {\begin{array}{c}
    \Gamma, x : A, p : \mathsf{Id}_A(a, x) \vdash P \type \\
    \Gamma \vdash d : P[a, \mathsf{refl}] \at{w} \\
    \Gamma \vdash e : \mathsf{Id}_A(a, b) \at{v}
  \end{array}}
\end{gather}

%=============================================================================
\section{Universes}
%=============================================================================

\begin{gather}
\infer[\textsc{Univ}]
  {\Gamma \vdash \Type_i \type}
  {\Gamma \ctx}
\\[1em]
\infer[\textsc{W-Univ}]
  {\Gamma \vdash \W : \Type_0 \at{\one}}
  {\Gamma \ctx}
\\[1em]
\infer[\textsc{El}]
  {\Gamma \vdash A \type}
  {\Gamma \vdash A : \Type_i \at{\one}}
\end{gather}

Types-as-terms have weight $\one$. The weight monoid lives in $\Type_0$.

%=============================================================================
\section{MLTT Embedding}
%=============================================================================

\begin{theorem}
Let $\W = \{\zero, \one\}$ with $\one \cdot \one = \one$, else $\zero$.

Then ProbTT judgments reduce to MLTT judgments:
\begin{itemize}
    \item $\Gamma \vdash a : A \at{\one}$ becomes $\Gamma \vdash a : A$
    \item $\Gamma \vdash a : A \at{\zero}$ is vacuously true (ignore)
    \item All rules with $w \cdot v$ compute: $\one \cdot \one = \one$, else $\zero$
\end{itemize}
\end{theorem}

\begin{proof}
By inspection. When weights are Boolean:
\begin{itemize}
    \item \textsc{Fun-Elim}: $\one \cdot \one = \one$, so certain function + certain argument = certain result
    \item \textsc{Pair-Intro}: $\one \cdot \one = \one$, so certain components = certain pair
    \item \textsc{Weaken}: $v \leq w$ with $v, w \in \{\zero, \one\}$ is trivial when $w = \one$
\end{itemize}
The rules become exactly MLTT.
\end{proof}

%=============================================================================
\section{What We Gain}
%=============================================================================

\subsection{Graded Reasoning}

With $\W = [0, 1]$:

\begin{itemize}
    \item $f : A \to B \at{0.9}$ means: $f$ produces valid output 90\% of the time
    \item $(a, b) : A \times B \at{0.8 \cdot 0.7}$ means: joint probability 0.56
    \item Composition multiplies: uncertain reasoning chains degrade gracefully
\end{itemize}

\subsection{Resource Interpretation}

With $\W = [0, \infty]$ (tropical):

\begin{itemize}
    \item $a : A \at{c}$ means: producing $a$ costs $c$
    \item Composition adds costs (since tropical $\cdot$ is $+$)
    \item This gives cost-aware type theory
\end{itemize}

\subsection{Possibility Interpretation}

With $\W = \{\zero, ?, \one\}$ (three-valued):

\begin{itemize}
    \item $\one$ = definitely true
    \item $?$ = possibly true
    \item $\zero$ = definitely false
    \item $? \cdot ? = ?$, $\one \cdot ? = ?$, etc.
\end{itemize}

%=============================================================================
\section{What About Addition?}
%=============================================================================

We deliberately omitted $+$. Where would we need it?

\subsection{Disjunction}

In probability: $P(A \vee B) = P(A) + P(B) - P(A \wedge B)$.

But this requires both $+$ AND $-$. Without them, we can still have sum types:

\begin{gather}
\infer[\textsc{Sum-Form}]
  {\Gamma \vdash A + B \type}
  {\Gamma \vdash A \type \quad \Gamma \vdash B \type}
\\[1em]
\infer[\textsc{Sum-InL}]
  {\Gamma \vdash \mathsf{inl}\,a : A + B \at{w}}
  {\Gamma \vdash a : A \at{w}}
\\[1em]
\infer[\textsc{Sum-InR}]
  {\Gamma \vdash \mathsf{inr}\,b : A + B \at{w}}
  {\Gamma \vdash b : B \at{w}}
\end{gather}

The weight of $\mathsf{inl}\,a$ is just $w$---no addition involved.

\subsection{Negation}

Classical negation uses complement: $\neg P = 1 - P$.

Without $-$, we define negation as:
\[
\neg A := A \to \mathbf{0}
\]

A proof of $\neg A$ with weight $w$ is a function that, given $a : A$, produces $\bot$ with weight $w$. This works without subtraction.

\subsection{When We'd Need Addition}

\begin{itemize}
    \item Mixing independent distributions
    \item Bayesian update (needs $+$ and $-$)
    \item Expected value (needs integration)
\end{itemize}

For these, extend $\W$ to a semiring. But the \textbf{core type theory works without}.

%=============================================================================
\section{Categorical Semantics}
%=============================================================================

\begin{definition}[Weighted Category]
A \emph{$\W$-weighted category} has:
\begin{itemize}
    \item Objects (types)
    \item For each $w : \W$, a set of morphisms $\mathrm{Hom}_w(A, B)$
    \item Identity: $\mathrm{id}_A \in \mathrm{Hom}_\one(A, A)$
    \item Composition: $f \in \mathrm{Hom}_w(A, B)$, $g \in \mathrm{Hom}_v(B, C)$ gives $g \circ f \in \mathrm{Hom}_{w \cdot v}(A, C)$
\end{itemize}
\end{definition}

\begin{theorem}[Soundness]
ProbTT has a sound interpretation in $\W$-weighted categories:
\begin{itemize}
    \item $\llbracket A \rrbracket$ = object
    \item $\llbracket \Gamma \vdash a : A \at{w} \rrbracket$ = morphism in $\mathrm{Hom}_w(\llbracket \Gamma \rrbracket, \llbracket A \rrbracket)$
    \item Weight multiplication = composition
\end{itemize}
\end{theorem}

When $\W = [0,1]$, this becomes a Markov category.

%=============================================================================
\section{Summary}
%=============================================================================

\textbf{Primitive}: Weight monoid $\W$ with $\zero, \one, \cdot, \leq$.

\textbf{Judgments}: $\Gamma \vdash a : A \at{w}$.

\textbf{Key rule}: Weights multiply in elimination.

\textbf{MLTT}: The $\W = \{\zero, \one\}$ case.

\textbf{Probability}: The $\W = [0,1]$ case.

\textbf{Markov structure}: Copy, delete, disintegration give conditioning.

\textbf{Chain rule}: $P(A|B) \cdot P(B) = P(A) \cdot P(B|A)$ (no division needed).

\textbf{No addition needed} for core type theory.

%=============================================================================
\section{Markov Structure and Bayes' Rule}
%=============================================================================

When $\W = [0,1]$, our weighted category becomes a \emph{Markov category}. This gives us conditioning and Bayes' rule without assuming real numbers as primitive.

\subsection{Markov Category Structure}

\begin{definition}[Copy and Delete]
A Markov category has additional structure:
\begin{itemize}
    \item \textbf{Copy}: $\Delta_A : A \to A \times A$ (diagonal)
    \item \textbf{Delete}: $!_A : A \to \mathbf{1}$ (terminal)
\end{itemize}
satisfying commutative comonoid laws.
\end{definition}

In ProbTT, these are:
\begin{gather}
\infer[\textsc{Copy}]
  {\Gamma \vdash \Delta\,a : A \times A \at{w}}
  {\Gamma \vdash a : A \at{w}}
\\[1em]
\infer[\textsc{Delete}]
  {\Gamma \vdash !\,a : \mathbf{1} \at{\one}}
  {\Gamma \vdash a : A \at{w}}
\end{gather}

Copy preserves weight (the same evidence supports both copies). Delete always succeeds with weight $\one$.

\subsection{Conditioning and Disintegration}

\begin{definition}[Joint and Marginal]
Given $f : \Gamma \to A \times B \at{w}$:
\begin{itemize}
    \item \textbf{Marginal}: $\pi_1 \circ f : \Gamma \to A$ and $\pi_2 \circ f : \Gamma \to B$
    \item \textbf{Joint}: The original $f$ encodes the joint distribution
\end{itemize}
\end{definition}

\begin{definition}[Disintegration]
A morphism $f : \Gamma \to A \times B \at{w}$ \emph{disintegrates} if there exist:
\begin{itemize}
    \item $g : \Gamma \to A \at{v}$ (marginal on $A$)
    \item $h : \Gamma \times A \to B \at{u}$ (conditional $B$ given $A$)
\end{itemize}
such that $f = \langle g, h \circ (\mathrm{id} \times g) \rangle$ with $w = v \cdot u$.
\end{definition}

\subsection{Chain Rule (Weak Bayes)}

\begin{theorem}[Chain Rule]
In a Markov category admitting disintegration, the chain rule holds:

Given joint $P(A, B)$ that disintegrates as:
\[
P(A, B) = P(A) \cdot P(B \mid A)
\]
there exists (by symmetry) another disintegration:
\[
P(A, B) = P(B) \cdot P(A \mid B)
\]
Therefore:
\[
P(A \mid B) \cdot P(B) = P(A) \cdot P(B \mid A)
\]
\end{theorem}

\begin{remark}
This is \textbf{Bayes' rule in multiplicative form}---no division required! The weight monoid needs only multiplication. Division would give the ``strong'' form $P(A \mid B) = P(A) \cdot P(B \mid A) / P(B)$, but we avoid it to stay minimal.
\end{remark}

\subsection{Conditioning in ProbTT}

Conditioning is expressed via the chain rule, not division:

\begin{gather}
\infer[\textsc{Chain}]
  {\Gamma \vdash (a, b) : A \times B \at{u \cdot v}}
  {\Gamma \vdash a : A \at{u} \quad \Gamma \vdash b \mid a : B \at{v}}
\end{gather}

Read: ``joint weight = marginal $\times$ conditional.'' The conditional $b \mid a$ is \emph{primitive}---we don't compute it by division, we assert it exists when the chain rule holds.

\begin{remark}[No Division Needed]
In the {0,1} case, conditioning is trivial: $P(B \mid A) = 1$ when $A \Rightarrow B$, else $0$.

For $\W = [0,1]$, conditioning exists but we express it multiplicatively. This keeps the weight monoid minimal (no division operation required).
\end{remark}

\begin{remark}[Graded Ex Falso]
What about $P(A \mid B)$ when $P(B) = 0$? This is the probabilistic \textbf{ex falso quodlibet}.

In classical logic: $\bot \to A$ (from absurdity, anything follows).

In ProbTT: when $b : B \at{\zero}$, we have for any $a$:
\[
(a, b) : A \times B \at{\zero \cdot w} = \at{\zero}
\]
The conditional $a \mid b$ can have \emph{any} weight---it's vacuously satisfied because the joint is already zero.

\textbf{Key insight}: Ex falso is the $w = 0$ limit of a \emph{continuous} phenomenon:
\begin{itemize}
    \item $w = \one$: Evidence fully constrains conclusions (classical implication)
    \item $w = \epsilon$ (small): Evidence weakly constrains conclusions
    \item $w = \zero$: Evidence imposes no constraint (ex falso)
\end{itemize}
Classical logic has only the endpoints; ProbTT has the entire spectrum.
\end{remark}

\subsection{What Markov Categories Give Us}

\begin{enumerate}
    \item \textbf{Copying}: $\Delta$ lets us use the same evidence twice
    \item \textbf{Marginalization}: Delete/project to ``forget'' components
    \item \textbf{Conditioning}: Disintegration extracts conditionals from joints
    \item \textbf{Chain rule}: $P(A, B) = P(A) \cdot P(B \mid A)$ (multiplicative Bayes)
    \item \textbf{Independence}: $P(A, B) = P(A) \cdot P(B)$ when conditional equals marginal
\end{enumerate}

All without assuming real numbers or division---just the weight monoid with multiplication.

%=============================================================================
\section{Open Problems}
%=============================================================================

\begin{enumerate}
    \item \textbf{Weight inference}: Can we infer $w$ from term structure?

    \item \textbf{Decidability}: Type checking is decidable if $\W$ is decidable.

    \item \textbf{Normalization}: Should follow from MLTT since weights don't affect reduction.

    \item \textbf{Dependent weights}: What if $w$ depends on $x$ in $(x : A) \to B \at{w(x)}$?

    \item \textbf{Continuous types}: Extension to measurable spaces?

    \item \textbf{Implementation}: Type checker for ProbTT?
\end{enumerate}

\end{document}
